\subsection{Correspondencia con los 28 indicadores del plan de negocios}
\label{ap:rubrica}
Esta tabla indica dónde se atiende cada criterio y qué falta comprobar. Conserva el orden de los 28 indicadores, incluidos los repetidos m y o. La calificación corresponde al jurado \parencite{rubrica_plan_negocios_2026}.
\begin{longtable}{p{4.0cm}p{5.6cm}p{4.2cm}}
\caption{Correspondencia del informe con la rúbrica suministrada}\\
\toprule Indicador de la rúbrica & Desarrollo en el informe & Aspecto pendiente \\
\midrule\endfirsthead
\toprule Indicador de la rúbrica & Desarrollo en el informe & Aspecto pendiente \\
\midrule\endhead
a. Explica con solidez qué hace único al negocio y porqué es atractivo. & Idea, antecedentes y definición del negocio: Conexión de gestión administrativa y deportiva, acceso móvil y acompañamiento. & La preferencia frente a competidores debe comprobarse en demostraciones. \\
b. Describe las actividades clave que la empresa implementa para ofrecer una propuesta de valor. & Canvas y estudio técnico de operación: Investigación, desarrollo, pruebas, incorporación y soporte con responsables y registros. & La frecuencia y capacidad de soporte se contrastarán con uso real. \\
c. Evidencia la identificación de los equipos y recursos necesarios para llevar a cabo las operaciones de la empresa. & Inventario operativo: Equipo físico, conectividad, infraestructura, trabajo y sus costos presupuestados. & Falta completar el inventario y las cotizaciones; revisar si Workspace se contó dos veces. \\
d. Describe de forma detallada el producto o servicio propuesto que brindan valor a los clientes. & Servicio ofrecido y módulos funcionales: Miembros, membresías, pagos, asistencia, rutinas y progreso, con capturas del prototipo. & Aislamiento multiempresa y personalización siguen como desarrollo posterior. \\
e. Detalla cómo los productos o servicios ofrecidos se diferencian de la competencia. & Alternativas actuales y decisión de compra: Comparación de papel, hojas de cálculo, mensajería y plataformas; foco en tareas concretas. & No se afirma exclusividad ni superioridad comprobada frente a todo competidor. \\
f. Emplea lenguaje técnico acorde con el nivel académico y el campo del negocio. & Arquitectura y conceptos financieros: Se explican SaaS, roles, MRR, margen, liquidez, costos y recuperación. & El equipo debe poder explicar estos conceptos durante la exposición. \\
g. Define los canales mediante los cuales hará llegar a los clientes la propuesta de valor. & Canvas y plan de promoción: Contacto y demostración para venta; acceso web para entrega; mensajería para soporte. & Falta probar qué canales consiguen contrataciones. \\
h. Caracteriza ampliamente el segmento de clientes (Necesidades - comportamientos - atributos). & Cliente, usuario y mercado alcanzable: Comprador y beneficiarios diferenciados; atributos, dispositivos y necesidades de la muestra. & Solo cinco responsables y diez miembros; muestra no probabilística. \\
i. Presenta los elementos diferenciadores que facilitan la decisión de compra del cliente. & Competencia y decisión de compra: Demostración, alcance explícito, precio y acompañamiento de incorporación. & Falta conocer la respuesta de posibles compradores. \\
j. Describe las demandas del segmento de clientes y el seguimiento para asegurar la calidad de los bienes o servicios ofrecidos. & Calidad y diseño del piloto: Seguimiento por tarea, duración, éxito, errores, incidentes y revisión semanal. & El piloto permitirá medir las mejoras; sigue pendiente. \\
k. Presenta las estrategias para el acercamiento al cliente, ya sea durante el proceso de atención o de servicio, & Calendario comercial: Contacto inicial, demostración, seguimiento y soporte con recursos del presupuesto. & El calendario es propuesto y requiere registro de ejecución. \\
l. Presenta datos sobre clientes, tendencias y oportunidades. & Antecedentes y estudio de mercado: Datos de encuesta y contexto documentado de digitalización; oportunidad de centralizar tareas. & No existe censo local ni estimación representativa de demanda. \\
m. Incluye los canales para la distribución del producto hasta el cliente y su promoción, incorporando el uso de nuevas tecnologías. & Promoción, distribución y permanencia: Página, contenido demostrativo, contacto directo y distribución digital del servicio. & El criterio o solicita lo mismo; ambos remiten a esta explicación. \\
n. Presenta la estructura de costos, gastos e ingresos. & Costos y proyección anual: Se distinguen costos fijos, variables, trabajo, instalación, desarrollo e ingresos mensuales. & Falta confirmar precios y completar los gastos pendientes. \\
o. Incluye los canales para la distribución del producto hasta el cliente y su promoción, incorporando el uso de nuevas tecnologías. & Promoción, distribución y permanencia: Acceso web, demostraciones y canales digitales de promoción y atención. & Repite el criterio m de la rúbrica. \\
p. Determina el mercado meta y su alineación con la propuesta de valor del negocio. & Comprador, usuario y mercado alcanzable: Gimnasios con necesidad administrativa, conectividad y responsable de decisión. & Diez prospectos y cinco altas son hipótesis, no mercado cuantificado. \\
q. Presenta el modelo de negocio integrando los nueve módulos básicos que definen la operación de la empresa, según lo expuesto en la fase Institucional. & Desarrollo del Canvas y apéndice visual: Nueve módulos: segmentos, valor, canales, relaciones, ingresos, recursos, actividades, alianzas y costos. & Alianzas comerciales y aceptación de tarifas siguen por validar. \\
r. Presenta el análisis de la viabilidad financiera incluyendo ingresos, costos, gastos y utilidades esperadas. & Análisis económico anual: Tablas de ingresos, gastos, caja y balance, con explicación del margen y la recuperación. & Resultados previstos antes de impuestos; falta comprobarlos durante la operación. \\
s. Presenta el análisis de las fuentes de financiamiento y su impacto en la permanencia del negocio. & Financiamiento y permanencia: Aporte propio y trabajo; comparación con reinversión, crédito y capital externo. & Sin préstamo ni inversión externa confirmados; depende de disponibilidad del equipo. \\
t. Evidencia una relación realista entre las proyecciones e indicadores financieros. & Relaciones financieras y sensibilidad: Relación entre clientes, ingresos y costos; cambios si baja el ingreso o sube el gasto. & Revisar el trabajo sin pago y los impuestos antes de operar. \\
u. Evidencia cumplimiento de la normativa vigente en apego a la forma jurídica, organizacional y seguridad social. & Forma jurídica y evidencia de formalización: Trámites por consultar y pasos para registrar el negocio. & Inscripciones y afiliaciones pendientes. \\
v. Detalla las alianzas estratégicas y aportes a su propuesta de valor. & Alianzas y continuidad: Tutor confirmado; gimnasio potencial con aportes y condiciones; proveedores diferenciados. & No hay alianza comercial documentada ni aceptación de piloto. \\
w. Presenta una organización clara y lógica del documento, en congruencia con la estructura dada en los lineamientos. & Orden del informe e índice: Base institucional más definición, mercado, técnica, organización, inversión y evaluación anual. & Verificar el PDF completo y su correspondencia con lineamientos. \\
x. Presenta el listado de referencias citadas en el documento, según formato APA vigente. & Citas, notas de fuente y referencias: Fuentes consultadas citadas en APA y cálculos propios identificados en tablas y figuras. & La proyección financiera es elaboración propia; sus supuestos necesitan comprobación. \\
y. Cumple con el formato establecido. & Portada y maquetación: Carta, márgenes de 2,54 cm, cuerpo de 12 puntos, interlineado y numeración. & Control visual de tablas, figuras y referencias en el PDF final. \\
z. Justifica de forma sólida la viabilidad y pertinencia del negocio. & Justificación y conclusión financiera: Necesidad observada, prototipo y escenario con margen de 21,44 por ciento antes de impuestos. & La viabilidad comercial está condicionada a piloto, capacidad y aceptación de precio. \\
aa. Identifica la situación de mercado, sociedad o industria que se aborda en la propuesta de negocio. & Situación del mercado e industria: Gestión dispersa y contexto de digitalización con fuentes y muestra identificadas. & No se extrapola al total de establecimientos de Pérez Zeledón. \\
bb. Argumenta la necesidad o problema que resuelve la propuesta de negocio. & Problema, encuesta y propuesta de valor: Dificultades de cobro, vigencia y seguimiento vinculadas con funciones concretas. & El efecto de la solución requiere comparación de tareas durante el piloto. \\
\bottomrule
\end{longtable}
\textit{Nota}. La tabla permite revisar los 28 indicadores y las tareas pendientes antes de la entrega.
